% ===========================================================================
% Annexe B — Synthèse de la campagne de validation.
% ===========================================================================
\chapter{Synthèse de la campagne de validation}
\label{app:validation_campaign}

Cette annexe consigne la campagne du \textbf{22 septembre 2026}. Elle doit être
reprise et redatée après correction des trois échecs analysés à la
section~\ref{sec:failure_analysis}.

\section{Conditions d'exécution}

\begin{table}[htbp]
  \centering
  \caption{Conditions d'exécution de la campagne}
  \label{tab:campaign_conditions}
  \begin{tabularx}{\textwidth}{L{4.0cm} X}
    \toprule
    \textbf{Élément} & \textbf{Valeur} \\
    \midrule
    Date de la campagne   & 22 septembre 2026 \\
    Base de données       & Base embarquée en local ; moteur cible avec extension vectorielle en intégration continue \\
    Langage serveur       & Python 3.14 \\
    Interface             & Angular 19.2, TypeScript 5.7 \\
    Portée                & Suite serveur complète et build de production de l'interface \\
    Hors portée           & Tests de bout en bout, de charge, de résilience et de sécurité dynamique \\
    \bottomrule
  \end{tabularx}
\end{table}

\section{Résultats bruts}

\begin{table}[htbp]
  \centering
  \caption{Résultats bruts de la campagne}
  \label{tab:campaign_raw}
  \begin{tabularx}{\textwidth}{L{6.0cm} X}
    \toprule
    \textbf{Mesure} & \textbf{Valeur relevée} \\
    \midrule
    Tests réussis                 & 594 \\
    Tests échoués                 & \textbf{3} \\
    Tests ignorés                 & 6 \\
    Durée de la suite             & 607,4 s \\
    Code de sortie de la suite    & 1 \\
    Durée du build de production  & 36,0 s \\
    Code de sortie du build       & 0 \\
    Bundle initial (brut)         & 1,25 Mo \\
    Bundle initial (transféré)    & 246,82 ko \\
    Avertissements de budget      & 4 \\
    \bottomrule
  \end{tabularx}
\end{table}

\section{Détail des échecs}

\begin{table}[htbp]
  \centering
  \caption{Détail des trois échecs et qualification}
  \label{tab:campaign_failures}
  \begin{tabularx}{\textwidth}{L{1.2cm} L{2.2cm} L{2.6cm} X}
    \toprule
    \textbf{Réf.} & \textbf{Attendu / obtenu} & \textbf{Domaine} & \textbf{Qualification} \\
    \midrule
    E-01 & 200 / 502 & Synchronisation d'agenda
         & Défaillance du dispositif de test, dépendante de l'environnement. Ne démontre ni défaut ni bon fonctionnement du produit. \\
    E-02 & 200 / 403 & Autorisation
         & Attente de test périmée. Le refus est conforme à la politique en vigueur. Pas un défaut du produit. \\
    E-03 & 403 / 200 & Autorisation
         & \textbf{Défaut réel.} Un rôle que la politique entend exclure atteint une action engageant financièrement l'entreprise. L'endpoint est gardé par une permission plus large que celle visée. \\
    \bottomrule
  \end{tabularx}
\end{table}

\section{Métriques du modèle salarial}

Relevées dans l'artefact daté du 10 septembre 2026, relu le 22 septembre 2026.

\begin{table}[htbp]
  \centering
  \caption{Métriques du modèle de prédiction salariale}
  \label{tab:campaign_model}
  \begin{tabularx}{\textwidth}{L{6.0cm} X}
    \toprule
    \textbf{Métrique} & \textbf{Valeur} \\
    \midrule
    Erreur absolue moyenne (test) & 402 TND \\
    Coefficient de détermination  & 0,891 \\
    Modèles comparés              & 8 \\
    Partitions d'évaluation       & 5 \\
    Critère de sélection          & Erreur de test, non d'entraînement \\
    \bottomrule
  \end{tabularx}
\end{table}

\section{Vocabulaire d'états employé}

\begin{table}[htbp]
  \centering
  \caption{Vocabulaire d'états de preuve}
  \label{tab:campaign_vocab}
  \begin{tabularx}{\textwidth}{L{4.6cm} X}
    \toprule
    \textbf{État} & \textbf{Signification stricte} \\
    \midrule
    Réussite capturée        & Une exécution datée existe et a réussi \\
    Échec capturé            & Une exécution datée existe et a échoué \\
    Implémentation confirmée & Le mécanisme existe dans le code ; aucune exécution n'a été capturée \\
    Partielle                & Une partie seulement du comportement a été exercée \\
    Planifié                 & Prévu, non exécuté \\
    Non vérifiée             & Aucune preuve n'existe \\
    Sans objet               & La validation ne s'applique pas \\
    \bottomrule
  \end{tabularx}
\end{table}

L'expression « réussite supposée » ne fait pas partie de ce vocabulaire, et
n'apparaît nulle part dans ce rapport.
